% ── Part IV, Chapter 4 ───────────────────────────────────────
% The Problem He Cannot Fix — limits of the embodied approach,
% what requires more than one person moving through a system
% Julius Caesar Scaliger, Pacific, late 21st century

\chapter{The Problem He Cannot Fix}

\strandmarker{Julius Caesar Scaliger}{Pacific, late 21st century}

\lettrine{T}{here} is a problem he has been circling for months.

Not the western grid, which resolved. Not the mapping
tool bug, which was fixed three weeks after he posted
the comment, the fix clean, the interaction between
the two functions now handled correctly, the wrong
outputs gone. Those problems had the quality of
problems that can be solved: bounded, locatable,
addressable by a specific intervention at a specific
point in the system.

This one does not have that quality.

He noticed it first in the kelp, which is where
he notices most things: a pattern of stress that
did not correspond to any of the conditions he
knows how to read. Not current, not temperature,
not the kind of biological pressure that comes
from competition or disease. Something else,
distributed across a section of the farm large
enough that he cannot see the whole of it from
any position he can reach, affecting the kelp
in a way that is consistent enough to be a
pattern and subtle enough that he is not certain
the pattern is real rather than a coincidence of
unrelated variations he is connecting because
connection is what his attention does when it
moves across things.

He has been trying to determine which.

\section{}

The difficulty is that he cannot get far enough
outside the system to see it.

This is a constraint he has not encountered in
the same form before. The western grid problem
was inside the system but visible from inside the
system: he could reach the affected area, observe
it directly, accumulate enough local evidence to
understand what was happening at that location.
The code bug was visible from outside the system:
he could read the code, trace the interaction,
identify the location of the failure without being
inside the system that was failing.

This problem is neither.

It is too large to observe locally and he does not
have access to the data that would allow him to
observe it globally. The sensors measure what the
sensors measure, and what the sensors measure is
being processed by systems he cannot read, and the
outputs of those systems are not available to him
in any form that would let him check whether what
he is seeing in the kelp is already known and
being addressed or is not known and should be
reported or is known and has been determined not
to require a response for reasons he would not
have access to even if he knew the determination
had been made.

He is, for the first time in a long time, working
at the edge of what a single person moving through
a system can find.

\section{}

He tries several approaches.

He maps the affected area as precisely as he can
from inside it, marking which lines show the stress
pattern and which do not, looking for a boundary
that might indicate a source. The boundary, if it
exists, is not where he can reach it. The pattern
extends past the range of his daily movement and
he cannot extend his range without equipment he
does not have and authorisation he would need to
request and a reason for requesting it that he
cannot give without already having the information
the extended range would provide.

He goes back to the logs.

The logs show him what the sensors show, which is
not the same as what he sees, and the gap between
them is itself information: whatever is producing
the pattern in the kelp is not producing a signal
that the sensors are measuring, which means either
the sensors are not measuring the right things or
the signal is below the sensors' resolution or the
signal is of a kind the sensors were not designed
to detect.

He has been in this position before, in smaller
versions of the same situation, and the resolution
has always been the same: the problem that a single
observer cannot fully locate requires more observers,
differently positioned, sharing what they find.

He does not have more observers.

He is Lord Erasmus on the ocean current mapping
forum and on three other trackers and in two
communities he reads more than he contributes to,
but none of these people are here, in this water,
able to see what he sees.

\section{}

He writes a description of what he has observed.

This is harder than the comment on the mapping
tool, which concerned a specific failure in a
specific system that he could point to precisely.
This concerns a pattern he is not certain is real,
observed through a method he cannot fully explain,
in a system too large for him to have observed
completely. He is describing the edge of his
knowledge, which is a different thing from
describing what he knows.

He writes it carefully, the way he writes anything
that matters, not drafting and revising but waiting
for each sentence to be complete before he begins it.

He describes what he has seen and where he has seen
it and over what period of time and with what degree
of confidence, which is not high. He describes the
approaches he has tried and what they did and did
not reveal. He describes the shape of the gap between
what he can observe and what he would need to observe
to determine whether the pattern is real.

He posts it in the one community where he thinks
someone might know how to read it.

Then he goes to sleep, because the farm will begin
again in the morning and he will move through it
and the pattern will either be more visible than
it was today or it will not be, and he cannot do
anything about which it will be by remaining awake.

\section{}

In the morning there is one response.

It is from someone he has never interacted with,
whose handle he does not recognise, who has been
working on a different oceanic farm two thousand
kilometres north and has seen something similar
and also does not know what it is.

He reads this for a long time.

Two observers, differently positioned, seeing the
same shape in the data. The shape is therefore
more likely to be real than coincidental, which
does not mean he knows what it is, only that it
is more worth the effort of finding out.

He writes back.

They begin, carefully, without forcing it, to
compare what they have each seen, building from
two partial observations toward something that
neither observation alone could have produced.
The process is slow. They are both working, both
moving through farms that require their attention,
both contributing to the shared picture in the
gaps between the demands of the day.

He does not know where it will go.

He knows it is the right process for the kind
of problem this is: too large for one person,
requiring more eyes differently placed, the
answer distributed across observers the way
the pattern is distributed across water, not
available from any single position but present
in the structure that the positions, taken
together, can begin to describe.

He puts the screen down and goes out to the farm.

The water is the colour it was yesterday, which
means the conditions have not changed, which
means the pattern is still there, patient,
indifferent to whether he has understood it yet.

He moves toward it anyway.

